% 表：假设扰动灵敏度分析
% 数据来源：figures/sensitivity_results.json, all_results.json, problem_3_results.json
\begin{table}[H]
  \centering
  \caption{关键假设扰动对达标时刻的影响}
  \label{tab:sensitivity}
  \small
  \begin{tabular}{llrrr}
    \toprule
    编号 & 扰动情景 & 基准 $t^*$ / h & 扰动 $t^*$ / h & 相对变化 \\
    \midrule
    H3 & 问题3 环境保持$\to$区间均值   & 57.262 & 57.567  & $+0.53\%$ \\
    H3 & 问题3 环境保持$\to$线性外推   & 57.262 & 57.115  & $-0.26\%$ \\
    H3 & 问题4 环境保持$\to$区间均值   & 50.898 & 51.166  & $+0.53\%$ \\
    H3 & 问题4 环境保持$\to$线性外推   & 50.898 & 50.782  & $-0.23\%$ \\
    H4 & 计入蒸发潜热 $L=2.4$ MJ$\cdot$kg$^{-1}$ & 57.262 & 62.646 & $+9.40\%$ \\
    H7 & 问题4 退化为固定域（忽略收缩） & 50.898 & 129.253 & $+153.95\%$ \\
    M3 & 界面平均 Kirchhoff$\to$算术    & 57.262 & 56.230  & $-1.80\%$ \\
    M3 & 界面平均 Kirchhoff$\to$调和    & 57.262 & 557.761 & $+874.05\%$ \\
    \bottomrule
  \end{tabular}

  \vspace{2pt}
  \parbox{0.94\linewidth}{\footnotesize
  注：环境外推假设（H3）影响不足 $0.6\%$，结论稳健。
  潜热（H4）单向延长干燥期，题面未给 $L$，故基准不含潜热汇并作为下界报告。
  忽略收缩（H7）使 $t^*$ 虚高 1.5 倍，证实动边界为问题4 的必要机制。
  调和平均在 $D$ 跨三个量级时被最小值锁死，产生 557.761 h 的非物理结果，
  故界面扩散系数采用 Kirchhoff 积分平均。}
\end{table}
